\section{Background and Motivation}
\label{sec:prelim}

\subsection{Asymmetric Quantization for KV Cache}

KV cache quantization reduces the storage footprint of attention keys and values by representing them in lower precision.
Following \kivi{}~\citep{kivi}, we adopt asymmetric (affine) quantization, where each quantized value $\hat{x}$ maps to its float16 representation as:
\begin{equation}
x = s \cdot (\hat{x} - z)
\label{eq:dequant}
\end{equation}
where $s$ is the per-channel (for keys) or per-token (for values) scale factor and $z$ is the zero-point.
For 2-bit quantization, $\hat{x} \in \{0,1,2,3\}$, requiring only 0.25 bytes per element — an $8\times$ reduction from float16.

\subsection{Decomposed Attention Dot Product}

The attention score between query $Q$ and key $K_i$ is computed as the dot product $Q \cdot K_i = \sum_{j=1}^{d} q_j k_{ij}$.
With asymmetric quantization, $q_j = s_Q(\hat{q}_j - z_Q)$ and $k_{ij} = s_K(\hat{k}_{ij} - z_K)$.
Substituting and expanding:
\begin{equation}
\begin{aligned}
Q \cdot K_i &= \sum_{j=1}^{d} s_Q(\hat{q}_j - z_Q) \cdot s_K(\hat{k}_{ij} - z_K) \\
&= s_Q s_K \left[\sum_{j=1}^{d} \hat{q}_j \hat{k}_{ij} - z_K\sum_{j=1}^{d} \hat{q}_j - z_Q\sum_{j=1}^{d} \hat{k}_{ij} + d \cdot z_Q z_K\right]
\end{aligned}
\label{eq:full_expansion}
\end{equation}

This decomposition reveals \textbf{four distinct operations}:
\begin{enumerate}
\item \textbf{Main dot product}: $\sum \hat{q}_j \hat{k}_{ij}$ — a standard GEMV on 2-bit data.
\item \textbf{Q-reduction}: $\sum \hat{q}_j$ — sum of query elements (computed once per query).
\item \textbf{K-reduction}: $\sum \hat{k}_{ij}$ — sum of key elements (per key vector).
\item \textbf{Constant term}: $d \cdot z_Q z_K$ — negligible.
\end{enumerate}

On GPU, these operations are typically fused into a single kernel.
However, the reduction terms require additional floating-point operations that compete for GPU compute units.
\textbf{Key observation}: operations (1)-(3) all operate on \emph{quantized 2-bit data}, which can be processed efficiently by simple near-memory compute units, while only the final scaling multiplication $s_Q s_K$ requires float16 precision.

\subsection{Near-Bank Processing-in-Memory}

Near-bank PIM places lightweight processing elements (PEs) inside each DRAM bank, directly adjacent to the memory arrays~\citep{seshadri2017ambit,mutlu2020processing}.
Each bank has:
\begin{itemize}
\item A \textbf{row buffer} ($R$ bytes): data read from DRAM arrays in a single activation, serving as the PE's working memory.
\item A \textbf{PE}: a simple multiply-accumulate unit processing $W$ bytes per cycle at frequency $f$.
\item \textbf{Local bandwidth}: $B$ GB/s from DRAM array to row buffer per bank.
\end{itemize}

% TODO Figure 3: Near-bank PIM microarchitecture — DRAM array → row buffer → PE

The total system aggregates $N$ such banks across multiple HBM stacks.
For a representative HBM3E configuration with 8 stacks, each with 256 PIM-enabled banks, $N = 2048$ banks operate in parallel.
Each bank processes its share of the data independently, with total latency determined by the slowest bank.

The key advantage of near-bank PIM for quantized attention is the \textbf{data locality}: 2-bit KV data stays within the bank and is processed by the local PE, avoiding the costly data movement to GPU.
The row-buffer pipeline naturally overlaps data movement with computation: while the PE computes on one row buffer's worth of data, the next row buffer fills from DRAM.

\subsection{Limitations of Existing Approaches}

\textbf{GPU-only execution}: \kivi{} demonstrates that asymmetric 2-bit quantization preserves accuracy, but the dequantization and reduction overhead on GPU limits throughput gains.
With batch size 64 and 16K context, GPU spends significant cycles on the extra reduction terms in Eq.~\ref{eq:full_expansion}.

\textbf{Logic-die PIM}: Prior work~\citep{duplex} places compute units on a separate logic die, requiring data to traverse TSVs from memory banks to logic die.
This introduces bandwidth bottlenecks and additional latency, making it unsuitable for the fine-grained, bank-local operations needed by asymmetric quantized attention.

\textbf{Direct 2-bit PIM compute}: While computing directly on 2-bit data in near-bank PIM is efficient, dequantizing first to float16 inside PIM (via $s \cdot (\hat{x} - z)$) incurs significant PE overhead.
Our experiments show this dequantization adds $3\times$ PE time, doubling total Q@K latency.
This motivates our approach of \emph{decomposing the asymmetric dot product} to keep as many operations as possible in the 2-bit domain.

% TODO Table 1: Comparison of compute substrates — GPU vs logic-die PIM vs near-bank PIM
% Columns: Compute throughput, memory bandwidth, 2-bit support, latency per GEMV
